%!TEX program = pdflatex
% CVPR/ICCV Conference Paper Template
% 适用于CVPR、ICCV、ECCV等计算机视觉会议

% 正确加载 IEEEtran 模板（CVPR官方使用）
\documentclass[10pt,twocolumn,letterpaper]{IEEEtran}

% 编码设置
\usepackage[utf8]{inputenc}
\usepackage[T1]{fontenc}

% 特殊字符支持
\usepackage{amsmath}
\usepackage{amssymb}
\usepackage{amsfonts}
\usepackage{amsthm}

% 图表支持
\usepackage{graphicx}
\usepackage[width=\linewidth]{figures/fig1.png}
\usepackage{subfig}
\usepackage{float}
\usepackage{booktabs}

% 表格
\usepackage{tabularx}
\usepackage{colortbl}
\usepackage{url}

% 算法/代码
\usepackage{algorithm}
\usepackage{algorithmicx}
\usepackage[plain]{algorithm2e}

% 参考文献
\usepackage{cite}

% 超链接
\usepackage[colorlinks=true, linkcolor=blue, citecolor=blue, urlcolor=blue]{hyperref}

% PDF元数据
\hypersetup{
    pdftitle={Paper Title},
    pdfauthor={Author Name},
    pdfkeywords={computer vision, deep learning, neural network},
}

% 自定义命令
\DeclareMathOperator*{\argmax}{argmax}
\DeclareMathOperator*{\argmin}{argmin}
\DeclareMathOperator*{\sign}{sign}
\providecommand{\IEEEauthorrefmark}[1]{$^{\ #1}$}

% 定理环境
\newtheorem{mydef}{Definition}
\newtheorem{mylemma}{Lemma}
\newtheorem{mythm}{Theorem}

% 页面设置
\input{settings.tex}

% Title
\title{Your Paper Title Here}

% Author
\author{
    \IEEEauthorblockN{Author Name$^{\dagger\ddagger}$}
    \IEEEauthorblockA{\textit{Institution Name}\\
        \textit{City, Country}\\
        email@example.com}
    \IEEEauthorblockN{Co-Author Name$^{\dagger}$}
    \IEEEauthorblockA{\textit{Another Institution}\\
        \textit{City, Country}\\
        coauthor@example.com}
    \thanks{$^{\dagger}$First affiliation. $^{‡}$Second affiliation. This work was supported by ...}
}

% 会议信息（CVPR不需要手动添加，会自动生成）
\conferenceinfo{CVPR}{2026}

\begin{document}

\maketitle

\begin{abstract}
Your abstract goes here. It should summarize the main contributions of your paper in a concise manner. CVPR abstracts are typically limited to 200 words. Include the problem statement, your approach, and the main results.
\end{abstract}

\begin{IEEEkeywords}
computer vision, deep learning, neural networks, object detection, image segmentation
\end{IEEEkeywords}

\input{intro.tex}

\input{related.tex}

\input{method.tex}

\input{experiments.tex}

\input{conclusion.tex}

\section*{Acknowledgment}
This work was supported by ...

% References
\bibliographystyle{ieeeettr}
\bibliography{references}

% 附录（可选）
\appendix
\section{Supplementary Material}
\label{sec:supp}
Additional theoretical proofs, more experimental results, or implementation details can be included here.

\end{document}
